% TOP_PROBLEM: {"id": 2845, "title": "Kirby Problem 3.47", "category_id": 7, "category": {"id": 7, "name": "topology", "display_name": "Topology", "description": "Properties preserved under continuous deformations.", "slug": "topology", "order_index": 7, "created_at": "2026-07-31T15:26:25.671Z"}, "status": "open"}
% FIRST_POSTED: null
% ENTRY_KIND: statement_only
% Imported from: batch12.tex; UnsolvedMath ID 2845
\documentclass[11pt]{article}
\usepackage[margin=25mm]{geometry}
\usepackage{fontspec}
\setmainfont{Latin Modern Roman}
\usepackage{amsmath,amssymb,mathtools}
\usepackage{unicode-math}
\setmathfont{Latin Modern Math}
\usepackage{xurl,hyperref}
\hypersetup{colorlinks=true,urlcolor=blue}
\setcounter{secnumdepth}{0}
\setlength{\parindent}{0pt}
\setlength{\parskip}{5pt}
\setlength{\emergencystretch}{3em}
\newcommand{\sref}[2]{\hyperlink{source:#1:#2}{[S#2]}}
\newcommand{\cataloguescope}[1]{\paragraph{Recorded target.}#1}
\begin{document}
% UnsolvedMath ID 2845; source catalogue code KP-3.47
\setcounter{section}{2844}
\section{Kirby Problem 3.47}\label{2845}
{\small\sffamily UnsolvedMath ID 2845 · Topology}
\subsection{Definitions and mathematical statement}

\paragraph{Shared notation.}$\mathbb N=\{1,2,\ldots\}$, and $\mathbb Z,\mathbb Q,\mathbb R,\mathbb C$ denote the integers, rationals, reals and complex numbers. Any other conventions are specified locally.


A contact form $\lambda$ on a closed $3$-manifold $Y$ is a smooth one-form with $\lambda\wedge d\lambda$ nowhere zero. Its Reeb vector field $R_\lambda$ is determined by $\lambda(R_\lambda)=1$ and $\iota_{R_\lambda}d\lambda=0$. A periodic orbit solves $\dot\gamma=R_\lambda\circ\gamma$ with some positive period. It is simple if it is not a multiple traversal of a shorter periodic orbit; simple orbits are counted by their images. On $\mathbb R^4=\mathbb C^2$ put $\lambda_0=\tfrac12\sum_{j=1}^2(x_j\,dy_j-y_j\,dx_j)$. The standard contact structure on $S^3$ is $\xi_{\rm std}=\ker(\lambda_0|_{S^3})$. A periodic Reeb orbit is elliptic if every eigenvalue of the linearized Poincar\'e return map on the transverse contact plane has modulus one; this includes degenerate unit eigenvalues. For a smooth bounded strictly convex domain $D\subset\mathbb R^4$, translate an interior point to the origin and use the contact form $\lambda_0|_{\partial D}$; strict convexity means that for any distinct $x,y\in\overline D$ and every $0<t<1$, $(1-t)x+ty\in D$. No positive-curvature condition on the boundary is added. \sref{2845}{3}
\paragraph{Statement recorded in the supplied source.}
(a) For every contact form $\lambda$ on $S^3$ defining $\xi_{\rm std}$, does its Reeb flow have an elliptic periodic orbit? (b) Does every Reeb flow obtained from the boundary of a smooth strictly convex domain in $\mathbb R^4$ have an elliptic periodic orbit? Both are universal existence questions, without genericity or pinching assumptions. \sref{2845}{2}
\subsection{Short English statement}
Does every standard-contact Reeb flow on the three-sphere have an elliptic closed orbit, and does this hold for every strictly convex hypersurface in four-dimensional Euclidean space?
\subsection{Sources}
\begin{description}
\item[\hypertarget{source:2845:1}{S1}] ULAM.AI and the UnsolvedMath Contributors, \emph{UnsolvedMath: A Curated Collection of Open Mathematics Problems}, record 2845 (KP-3.47). CC BY 4.0. \url{https://huggingface.co/datasets/ulamai/UnsolvedMath}.
\item[\hypertarget{source:2845:2}{S2}] R. İnanç Baykur, Robion C. Kirby and Daniel Ruberman (editors), \emph{K3: A New Problem List in Low-Dimensional Topology}, Mathematical Surveys and Monographs 295 (2026), Problem 3.47, p. 164 and following remarks; original author version. \url{https://math.berkeley.edu/sites/default/files/surv-295-ruberman-watermarked-author-pdf.pdf}.
\item[\hypertarget{source:2845:3}{S3}] Miguel Abreu and Leonardo Macarini, \emph{Dynamical convexity and elliptic periodic orbits for Reeb flows}, Mathematische Annalen 369 (2017), 331--386; arXiv:1411.2543v3, introduction, p. 1 and footnote 1. \url{https://arxiv.org/abs/1411.2543v3}.
\end{description}
\label{end:2845}
\end{document}
